\documentclass[11pt,a4paper,oneside,openany]{book}

\usepackage[utf8]{inputenc}
\usepackage[T1]{fontenc}
\usepackage[english]{babel}
\usepackage{csquotes}
\usepackage[charter]{mathdesign}
\usepackage{amsmath}
\usepackage{lipsum}
\usepackage{amsthm}
\usepackage{mathtools}
\usepackage{physics}
\usepackage{graphicx}
\usepackage{enumitem}
\usepackage{tikz-cd}
\usepackage[margin=2cm]{geometry}
\usepackage[backend=biber,style=numeric,sorting=none]{biblatex}
\addbibresource{bib.bib}
\usepackage[colorlinks=true,linkcolor=black,citecolor=black,urlcolor=black]{hyperref}

\newtheorem{theorem}{Theorem}[chapter]
\newtheorem{definition}{Definition}[section]
\newtheorem{lemma}[definition]{Lemma}
\newtheorem{proposition}[definition]{Proposition}
\newtheorem{corollary}[definition]{Corollary}
\theoremstyle{remark}
\newtheorem{remark}[definition]{Remark}
\newtheorem{example}[definition]{Example}

\let\mathcal\relax
\newcommand{\mathcal}[1]{\text{\usefont{OMS}{cmsy}{m}{n}#1}}
\DeclareMathOperator{\Spec}{Spec}
\DeclareMathOperator{\Ad}{Ad}
\DeclareMathOperator{\ad}{ad}
\newcommand{\poisson}[2]{\{#1,#2\}}
% ============================================================
% Comandi Matematici
% ============================================================

% --- Insiemi numerici ---
\newcommand{\R}{\mathbb{R}}
\newcommand{\C}{\mathbb{C}}
\newcommand{\N}{\mathbb{N}}
\newcommand{\K}{\mathbb{K}}
\newcommand{\A}{\mathfrak{A}}
\newcommand{\B}{\mathcal{B}}

% --- Spazio di Hilbert e spazio proiettivo ---
% ATTENZIONE: \H è già un comando di LaTeX (accento ungherese, es. \H{o}=ő).
% Il renewcommand lo sovrascrive: nella tesi non ti servirà mai quell'accento,
% quindi è sicuro, ma se preferisci non toccare comandi di sistema usa \Hil
% al posto di \H ovunque sotto.
\renewcommand{\H}{\mathcal{H}}
\newcommand{\PH}{\mathbb{P}\H}                    % P H, spazio proiettivo

% --- Algebre C* ---
\newcommand{\Cstar}{C^*}                          % da usare in math mode: $\Cstar$
\newcommand{\Csalg}{$C^*$-algebra}                % da usare nel testo: "è una \Csalg\ commutativa"
\newcommand{\Csalgs}{$C^*$-algebre}               % plurale
\newcommand{\BH}{B(\H)}                           % operatori limitati su H
\newcommand{\Cz}[1]{C_0(#1)}                      % C_0(M)
\newcommand{\Cinf}[1]{C^\infty(#1)}               % C^infty(M)
\newcommand{\Cinfc}[1]{C^\infty_c(#1)}            % C^infty_c(M)
\newcommand{\Mat}[2]{M_{#1}(#2)}                  % M_n(C), es. \Mat{2}{\C}

% --- Algebre di Lie (fraktur generico) ---
\newcommand{\mf}[1]{\mathfrak{#1}}                % \mf{g}, \mf{su}(2), \mf{u}(n), \mf{X}(M)

\newcommand{\qcomm}[2]{\dfrac{1}{i\hbar}[#1,#2]}  % (1/ih)[A,B]  -- condizione di Dirac, Qubit IV

% --- Notazione bra-ket ---
\newcommand{\melem}[3]{\langle#1|#2|#3\rangle}    % elemento di matrice <psi|A|psi>
\newcommand{\expect}[2]{\langle#1\rangle_{#2}}    % <A>_psi

% --- Campo vettoriale hamiltoniano ---
\newcommand{\Ham}[1]{X_{#1}}                      % \Ham{f}, \Ham{f_H}, \Ham{f_A}

% --- Simboli e operatori vari ---
\newcommand{\Id}{\mathbb{1}}
\renewcommand{\Re}{\mathrm{Re}}                   % la tesi usa Re/Im in tondo, non ℜ/ℑ
\renewcommand{\Im}{\mathrm{Im}}
\newcommand{\LC}{\varepsilon_{ijk}}               % simbolo di Levi-Civita (indici fissi ijk)
\DeclareMathOperator{\Ran}{Ran}
\DeclareMathOperator{\Ker}{Ker}
\DeclareMathOperator{\Span}{span}
\DeclareMathOperator{\sing}{sing}
\setcounter{chapter}{1}
\setcounter{section}{6}
\begin{document}
	\section{Classical Mechanics as a Commutative $\mathbf{C^*}$-algebra}
	The standard description of a classical system establishes geometry as the primary source of the whole framework. The Gelfand theory of commutative $C^*$-algebras introduces a change of perspective for the foundation of Hamiltonian mechanics. By exploiting the Gelfand-Naimark representation theorem one can reverse this relation and construct the theory starting from the abstract algebra of observables. However, this framework reconstructs only a Hausdorff topological space through the Gelfand spectrum of the algebra. In fact, the full picture will be recovered adding missing structure to enrich the topological space with properties such as differentiability. This approach dictates that the algebra of physical observables can provide a self-contained axiomatic basis for the formulation of classical mechanics.
	
	\subsection*{Basic topics in $\mathbf{C^*}$-algebras theory}
	To reconstruct classical theory algebraically, we first need to present $C^*$-algebras. This abstract mathematical structure recurs in both classical and quantum descriptions of observables. Specifically, we adopt this framework because involution naturally models the complex conjugation or the adjoint of physical quantities, while the underlying metric completeness ensures that the approximation of observables is rigorously well-defined. Here, we will present only the basic definitions to introduce the reader to the theory of commutative $C^*$-algebras following the approach of Brattelli and Robinson \cite[Chapter 2]{bratteliOperatorAlgebrasQuantum}.
	
	\noindent We begin by introducing the notion of an algebra, the basic structure on which the involution defining a $^*$-algebra will be built.
	
	\begin{definition}[Algebra]\label{def:algebra}
		An algebra $\A$ is a $\C$-vector space equipped with an associative and distributive multiplication law which is compatible with the other vector space operations.
	\end{definition}
	
	\begin{definition}[$^*$-algebra]\label{def:star-alg}
		Let $\A$ be an algebra as in Definition \ref{def:algebra}. An involution on $\A$ is a map associating an element $A \mapsto A^*$ satisfying three distinct conditions:
		\begin{itemize}
			\item $A^{**} = A$,
			\item $(AB)^* = B^* A^*$,
			\item $(\alpha A + \beta B)^* = \overline{\alpha} A^* + \overline{\beta} B^*$.
		\end{itemize}
		An algebra equipped with an involution map is called a $^*$-algebra. 
	\end{definition}
	
	Within a $^*$-algebra it is natural to isolate subsets that remain stable under the involution, since they will later correspond to physically meaningful collections of observables, as we will present in Definition \ref{def:smooth-subalgebra}. Therefore, we introduce the following notion of a selfadjoint subset.
	
	\begin{definition}[Selfadjoint subset]\label{def:selfadj-subset}
		Let $\A$ be a $^*$-algebra in the sense of Definition \ref{def:star-alg}. A subset $\mathfrak{B} \subset \A$ is called \emph{selfadjoint} if $A \in \mathfrak{B}$ implies that $A^* \in \mathfrak{B}$.
	\end{definition}
	
	\noindent Some particular elements of this subset are the adjoint of themselves. As we will see later in Definition \ref{def:obs-alg}, selfadjoint elements have a particular meaning in algebraic classical mechanics.
	
	\begin{definition}[Selfadjoint element]\label{def:selfadj-element}
		Let $\A$ be a $^*$-algebra in the sense of Definition \ref{def:star-alg}. An element $A \in \A$ is selfadjoint if $A = A^*$. The set of all the selfadjoint elements will be referred as \emph{selfadjoint sector}.
	\end{definition}
	
	\noindent Now, we endow the algebra with an additional structure that measures the length of its elements, called \emph{a norm}. This operation will let us control convergence and continuity within the algebra.
	
	\begin{definition}[Normed algebra]\label{def:norm-alg}
		An algebra $\A$, as in Definition \ref{def:algebra}, endowed with a map $A \in \A \mapsto \|A\| \in \R$ satisfying the requirements:
		\begin{itemize}
			\item $\|A\| \geq 0$ and $\|A\| = 0$ if, and only if, $A = 0$,
			\item $\|\alpha A\| = |\alpha| \|A\|$,
			\item $\|A + B\| \leq \|A\| + \|B\| \quad$ the triangle inequality,
			\item $\|AB\| \leq \|A\| \|B\| \quad $ the product inequality,
		\end{itemize}
        is called a normed algebra, while the map is called a norm.
	\end{definition}
	
	\noindent Normed algebras are relevant objects inside our construction, since the norm endows each of them with a notion of a particular topology called \emph{uniform topology}.
	
	\begin{definition}[Uniform topology]\label{def:uni-top}
		Let $\A$ be a normed algebra as in Definition \ref{def:norm-alg}. The norm defines a metric topology on $\A$ which is referred to as the uniform topology. The neighborhoods of an element $A \in \A$ are given by $$\mathscr{U}_{A, \varepsilon} = \{B \in \A \ | \ \|B - A\| < \varepsilon\},$$ where $\varepsilon \in \R^+$. 
	\end{definition}

	\noindent Analogously, endowing a vector space with a norm makes it a \emph{normed vector space}, and the uniform topology just introduced allows us to speak of convergent sequences within it. However, in a generic normed space, such sequences are not guaranteed to converge to an element within the space. Imposing the requirement of completeness leads to the concept of a \emph{Banach space}.

    \begin{definition}[Banach Space on $\C$]\label{def:banach-space}
    A \emph{Banach space} is a normed $\C$-vector space $(X, \|\cdot\|)$ which is complete to the metric induced by its norm.
    \end{definition}
    
    \noindent A Banach space can be provided with a multiplication law and an involution creating a Banach $^*$-algebra.

    \begin{definition}[Banach algebra]\label{def:banach-alg}
    A Banach algebra is a normed algebra $\mathfrak{B}$ over a $\C$-vector space, as in Definition \ref{def:norm-alg}, which is complete with respect to the metric induced by its norm.
    \end{definition}
    
    \noindent Since physical observables also carry an involution, we further require compatibility between this operation and the norm.
    
    \begin{definition}[Banach $^*$-algebra]\label{def:banach-star-alg}
    A Banach $^*$-algebra is a Banach algebra $\A$, as in Definition \ref{def:banach-alg}, equipped with an involution, as in Definition \ref{def:star-alg}, such that the norm satisfies the isometric property $$\|A\| = \|A^*\|, \qquad \forall A \in \A.$$
    \end{definition}
	
	Completeness assures that every Cauchy sequence in the space converges to an element thereof, thus guaranteeing that approximations of observables are well-defined. In a Banach $^*$-algebra, the norm and involution are constrained by $\|A\|=\|A^*\|$. Imposing the stronger identity $\|A^*A\|=\|A\|^2$ makes the norm compatible with the algebraic structure and defines a $C^*$-algebra.
	
	\begin{definition}[$C^*$-algebra]\label{def:c-star-algebra}
		A $C^*$-algebra is a Banach *-algebra $\A$, as in Definition \ref{def:banach-star-alg}, with the property $\|A^* A\| = \|A\|^2$ $\forall A \in \A$.
	\end{definition}
	
	C*-algebras constitute a key framework in mathematical physics where algebraic structures are canonically endowed with a complete metric topology. The abstraction of bounded operators leads, both in classical and in quantum physics, to the study of Banach algebras and their involutions. This approach allows the characterization of states and observables in both theories. A key class of $C^*$-algebras, particularly relevant for this reconstruction, consists of those possessing an identity element.
	
	\begin{definition}[Unital $C^*$-algebra]\label{def:unital-c-star-algebra}
		A unital C*-algebra is a C*-algebra $\A$, as in Definition \ref{def:c-star-algebra}, containing an identity element $\Id$ such that $A = \Id A = A \Id$ for all $A \in \A$.
	\end{definition}
	
	The identity introduced above plays the role of a distinguished element acting trivially on every other one. Before relying on this notion, it is worth verifying that, when it exists, it is uniquely determined by the algebra.
	
	\begin{proposition}
		A C*-algebra $\A$ can possess at most one identity element.
	\end{proposition}
	
	\begin{proof}
		If a second identity element $\Id'$ exists it must satisfy $\Id' = \Id \Id' = \Id$.
	\end{proof}
	
	\subsection*{Algebraic description of Hamiltonian Mechanics}
	We now apply the abstract framework developed above to the specific case of Hamiltonian mechanics. In the classical setting, physical intuition suggests starting from the phase space and only afterward considering functions defined on it. Our goal will be to indicate that the reverse order is equally legitimate. We now carry out this reconstruction explicitly, showing how the familiar objects of classical mechanics, observables, states, and dynamics, can be redescribed in the framework of commutative $C^*$-algebras.
	 
    \smallskip
    \paragraph{Observables.}
	
	We begin by algebraically characterizing the observables used in Hamiltonian mechanics. Let $\Gamma$ be a generic compact phase space, in the sense of Definition \ref{def:phase_space}, for simplicity, we assume that there exist global coordinates $(q,p)$ on $\Gamma$. Let $\mathcal{P}$ be the algebra of \emph{polynomial observables} generated by the phase space coordinates $q$ and $p$. These, operationally, correspond to the measurement of position or momentum of a physical system. Given that the coordinates have different values for different points on $\Gamma$, they separate points. $\mathcal{P}$ also includes constant functions, consequently it fulfills the hypotheses of the Stone-Weierstrass Theorem \cite[Theorem IV.9]{reedMethodsModernMathematical2003}. This result establishes that $\mathcal{P}$ is dense within $C^0(\Gamma; \R)$ under the uniform norm topology that, in this space, is defined analytically as:  \begin{equation}
    \| f \|_\infty = \sup_{p \in \Gamma} |f(p)| \qquad \forall f \in \mathcal{P}.\label{eq:sup-norm} 
    \end{equation} Taking the closure with respect to this norm exactly recovers the full space $\overline{\mathcal{P}}^{\| \cdot \|_\infty} = C^0(\Gamma; \R)$.  
	
    In a general \emph{classical scenario}, the phase space $\Gamma$ is abstracted to a compact Hausdorff space $X$, replacing $C^0(\Gamma; \R)$ with the classical continuous algebra in Definition \ref{def:gel-alg}. Physical observables are then directly recovered as the selfadjoint elements invariant under complex conjugation via Definition \ref{def:obs-alg}.
	
	\begin{definition}[Classical Continuous Algebra]\label{def:gel-alg}
		Let $X$ be a compact Hausdorff space. The \emph{classical continuous algebra} is the commutative unital $C^*$-algebra $\A = C^0(X; \C)$, as in Definition \ref{def:unital-c-star-algebra}.
	\end{definition}

    \noindent In the case above, $X=\Gamma$ and $\A = C^0(\Gamma)$.

    \begin{remark}
        The space $C^0(X; \C)$ is a commutative unital $C^*$-algebra equipped with a pointwise multiplication $(fh)(x)=f(x)h(x)$, the identity element $\Id(x)=1$, the involution $f^*(x)=\overline{f(x)}$ and the supremum norm in Equation \eqref{eq:sup-norm}. While the $C^*$-algebra condition follows from $\|f^*f\|_\infty = \|f\|_\infty^2$.
    \end{remark}
    
    This construction establishes that the algebra of classical observables can be modeled by a commutative unital $C^*$-algebra. Later in this section, this will constitute a fundamental component for the algebraic reconstruction of classical mechanics, as seen in Definition \ref{def:cas}. Since complex-valued elements lack a direct physical interpretation, physical \emph{observables} are strictly identified with the selfadjoint elements of the algebra.
	
	\begin{definition}[Classical Observable Algebra]\label{def:obs-alg}
		Let $\A$ be a commutative unital $C^*$-algebra, as in Definition \ref{def:unital-c-star-algebra}. The classical observable algebra $\A_{sa}$ is the selfadjoint sector, as seen in Definition \ref{def:selfadj-element}.
	\end{definition}
	
	\noindent A more detailed discussion about the physical quantities really measurable will be done in Subsection \ref{sec:interpret-phys}.
	
	\medskip
    \paragraph{States.}
	Having analysed the algebraic nature of observables, we now turn to states. From an operational point of view, if we prepare the system in a particular state $\omega$ and we measure an observable $n$ times, then we can define its mean value as
	\[
		< f >^{(\omega)}_n \ := \frac{1}{n}\sum_{i=1}^n m^{(\omega)}_i(f) \qquad \forall f \in \A_{sa},
	\]
	where $m^{(\omega)}_i(f)$ represent the $i$-th measurement and $\A_{sa}$ is the classical observable algebra in Definition \ref{def:obs-alg}. The limit of this expression, assuming that it converges experimentally, can be defined dually as
	\[
		\omega(f) := \lim_{n \to \infty} < f > ^{(\omega)}_n \qquad \forall f \in \A_{sa}.
	\]
	In this picture, a state $\omega$, acting on observables, yields the average of the measurements. This definition can however be extended to the whole unital commutative $C^*$-algebra $\A$. The state $\omega$ becomes a linear functional that must preserve positivity $\omega(f^*f) \geq 0$ for all $f \in \A$. In addition, it can be shown that positivity implies the Cauchy-Schwartz Inequality, as proved by \cite[Lemma 2.3.10]{bratteliOperatorAlgebrasQuantum}, namely
	\begin{equation}
	    |\omega(A^*B)| \leq (\omega(A^*A) \ \omega(B^*B))^{1/2} \quad \forall A,B \in \A.\label{eq:cs-ineq}
	\end{equation}
	Therefore, if $\A$ is unital, we can set $B = \Id$ in Equation \eqref{eq:cs-ineq} to obtain $|\omega(A^*)|^2 \leq \omega(A^*A) \, \omega(\Id)$ for all $A \in \A$. If $\omega(\Id)$ were zero, this inequality would imply $\omega(A^*) = 0$ for every $A \in \A$, thus forcing the functional $\omega$ to be identically zero. Consequently, for any non-trivial positive linear functional, it must strictly hold that $\omega(\Id) > 0$. This ensures that one can always normalize such a functional to define a proper state via the map $$\omega \mapsto \omega_{norm} = \omega(\Id)^{-1} \omega.$$
 	
 	\noindent We can formalize this concept in the following definition.
 	
 	\begin{definition}[State on a $C^*$-Algebra]\label{def:state}
 		Let $\A$ be a unital $C^*$-algebra as in Definition \ref{def:unital-c-star-algebra}. A \emph{state} on $\A$ is a linear functional $\omega:\A\to\mathbb{C}$ such that
 		\begin{enumerate}
 			\item $\omega(A^*A)\ge 0$ for all $A\in\A$ $\quad$ positivity,
 			\item $\omega(\Id)=1$  $\quad$ normalization.
 		\end{enumerate}
 		The set of all states is denoted by $\mathcal{S}(\A)$.
 	\end{definition}
 	
 	States on a $C^*$-algebra are not isolated objects, since a convex combination of two physically admissible preparations of a system is itself a physically admissible preparation. Thus, we can analyse the properties of this space.
 	
 	\begin{proposition}[Convexity of the State Space]\label{prop:state-convex}
 		Let $\A$ be a unital $C^*$-algebra as in Definition \ref{def:unital-c-star-algebra}. The set $\mathcal{S}(\A)$ of all states on $\A$ is convex.
 	\end{proposition}
 	
 	\begin{proof}
 		Let $\omega_1,\omega_2 \in \mathcal{S}(\A)$ and let $t \in [0,1]$. Define $\omega = t\omega_1 + (1-t)\omega_2$ which is linear on $\A$. Positivity follows because for any $A\in\A$,$\omega(A^*A) = t\,\omega_1(A^*A) + (1-t)\,\omega_2(A^*A) \ge 0.$
	 	normalization holds since $\omega(\Id) = t\,\omega_1(\Id) + (1-t)\,\omega_2(\Id) = t + (1-t) = 1$. Thus $\omega \in \mathcal{S}(\A)$ and $\mathcal{S}(\A)$ is convex.
 	\end{proof}

    \paragraph{Probabilistic interpretation of states.}
 	
 	One can ask which is the usual physical interpretation of such an abstract object. The solution comes from considering the classical scenario given by $X$, a compact Hausdorff space and $C^0(X)$ the associated classical continuous algebra as in Definition \ref{def:gel-alg}. It is proved by \cite[Proposition 2.3.11]{bratteliOperatorAlgebrasQuantum} that a state $\omega \in \mathcal{S}(C^0(X))$ is continuous. Therefore, by the Riesz-Markov Representation Theorem \cite[Theorem IV.14]{reedMethodsModernMathematical2003}, it defines a unique Borel measure $\mu_\omega$ on $X$ such that
 	\begin{equation}
 		\omega(f) = \int_X f \, d\mu_\omega, \quad \mu_\omega(X) = \omega(\Id) = 1.
 	\end{equation}
 	Thus, we can interpret the state as a probability distribution on the state space and observables as random variables. This probabilistic interpretation represents the core statistical mechanics.
 	
 	In this description, Hamiltonian mechanics is recovered as a special case. In fact, we can take $\mu_{P_0}$ as the probability measure concentrated on the point $P_0 = \{q_0, p_0\}$ of a compact phase space $\Gamma$, in the sense of Definition \ref{def:phase_space}. Then $\mu_{P_0}$ is a Dirac measure and the corresponding state $\omega_{P_0}$ is
 	\begin{equation}
 		\omega_{P_0}(f) = \int_\Gamma f \, d\mu_{P_0} = f(P_0).
 	\end{equation}
 	Such states are also called \textit{pure states} since they cannot be written as convex non-trivial linear combinations of other states. Measurements on pure states always yield the same result $f_\omega = \omega(f)$, so there is no dispersion and, probabilistically speaking, the variance of that state vanishes. We can formalize these concepts generally in the following way.
 	
 	\begin{definition}[Pure State]\label{def:pure-state}
 		Let $\A$ be a commutative unital $C^*$-algebra. A state $\omega\in\mathcal{S}(\A)$, as in Definition \ref{def:state}, is called \emph{pure} if it is an extreme point of the convex set $\mathcal{S}(\A)$. That is, whenever $\omega = t\omega_1+(1-t)\omega_2$ with $0<t<1$ and $\omega_1,\omega_2\in\mathcal{S}(\A)$, then $\omega_1=\omega_2=\omega$.
 	\end{definition}
 	
 	\noindent Now, we extend the statistical definition of variance to our algebraic framework. This notion will provide the precise sense in which pure states correspond to dispersion-free measurements.
 	
 	\begin{definition}[Variance of a State]\label{def:variance}
 		Let $\A$ be a unital $C^*$-algebra, as in Definition \ref{def:unital-c-star-algebra}, and let $\omega \in \mathcal{S}(\A)$ be a state, defined in Definition \ref{def:state}. The \emph{variance} of an element $A \in \A$ in the state $\omega$ is defined as
 		\[
 		\operatorname{Var}_\omega(A) := \omega(A^* A) - \omega(A^*) \omega(A).
 		\]
 	\end{definition}
 	
 	The quantity is non-negative by the Cauchy--Schwarz inequality \cite[Lemma 2.3.10]{bratteliOperatorAlgebrasQuantum}. In particular, since their measurements are dispersion-free, pure states show a strong form of this property.
 	
 	\begin{proposition}[Zero Variance for Pure States]\label{prop:pure-zero-var}
 	Let $\A$ be a commutative unital $C^*$-algebra as in Definition \ref{def:unital-c-star-algebra} and let $\omega \in \mathcal{S}(\A)$ be a pure state, shown in Definition \ref{def:pure-state}. Then
 	\[
 		\operatorname{Var}_\omega(A) = 0 \qquad \forall A \in \A.
 	\]
 	\end{proposition}
	\begin{proof}
		By a standard result for commutative $C^*$-algebras, every pure state is multiplicative, that is, $\omega(AB)=\omega(A)\omega(B)$ for all $A,B\in\A$, see \cite[Proposition 2.3.27]{bratteliOperatorAlgebrasQuantum}. In particular, $\omega(A^*A)=\omega(A^*)\omega(A)$. Substituting into Definition \ref{def:variance} yields $\operatorname{Var}_\omega(A)=0$ for every $A\in\A$.
	\end{proof}
	
 	\medskip
    \paragraph{Duality between states and observables.}
    
 	To conclude the excursus on observables and states, we establish an intrinsic duality between these two notions in classical scenarios. On one hand, given $\A$ the classical continuous $C^*$-algebra, defined in Definition \ref{def:gel-alg}, every state $\omega \in \mathcal{S}(\A)$ induces an evaluation. When that map is restricted to its classical observable algebra $\A_{sa}$ introduced in Definition \ref{def:obs-alg} it yields strictly real values. This property is mathematically guaranteed by the positivity of the state, ensuring that $A = A^* \in \A_{sa}$ implies $\omega(A) \in \mathbb{R}$. Consequently, because every element $A \in \A$ admits a unique decomposition into selfadjoint components as $A = A_1 + iA_2$, with $A_1, A_2 \in \A_{sa}$, the action of a state on the entire complex $C^*$-algebra is completely and uniquely determined by its values on the physical classical observables. 
 	
 	Conversely, the uniqueness of a state given its observable values relies on the topological properties of the space $X$ associated with the classical continuous algebra $\A$. Since $X$ is compact and Hausdorff, it satisfies the hypothesis of Urysohn's lemma \cite[Theorem IV.7]{reedMethodsModernMathematical2003}. It states that for any two distinct points $p$ and $q$ in $X$, there exists a strictly real-valued function $f \in C^0(X; \mathbb{R}) \cong \A_{sa}$ such that $f(p)=0$ and $f(q)=1$. In terms of pure states, this yields $\omega_p(f)=0$ and $\omega_q(f)=1$. Thus, the classical observable algebra $\A_{sa}$ separates points on $X$. 
 	
 	Combining these findings yields a duality relation. A physical state is completely characterized by the collection of its values on the selfadjoint observables, and the algebra of such observables is structurally sufficient to reconstruct the points of the underlying state space.
 	\medskip
	\paragraph{Time evolution.}
 	Having established observables and states, we now algebraically investigate dynamics in Hamiltonian mechanics. Let $\Gamma$ be a compact phase space, as in Definition \ref{def:phase_space}. In this framework, time evolution of the system $\{q, p\} \to \{q_t, p_t\}$ is continuous in time $t \in \R$, with a continuous dependence on the initial data at time $t_0$. Therefore it defines a one-parameter family of continuous invertible maps $\alpha_{t_0, t}$ of $C^0(\Gamma; \R)$ into itself. These maps preserve all the algebraic relations, including the $^*$-operation. Such maps are called $^*$-automorphisms. Thus, the dynamics is encoded by a one-parameter group of automorphisms of the observable algebra.
 	
 	\begin{definition}[Time Evolution]\label{def:time-evolution}
 		Let $\A$ be a unital $C^*$-algebra as in Definition \ref{def:unital-c-star-algebra}. A \emph{time evolution} is a strongly continuous one-parameter group of \emph{$^*$-automorphisms} $\alpha_t:\A\to\A$, $t\in\mathbb{R}$, such that
 		\begin{enumerate}
 			\item $\alpha_0=\mathrm{id}$,
 			\item $\alpha_{s+t}=\alpha_s\circ\alpha_t$ for all $s,t\in\mathbb{R}$,
 			\item each $\alpha_t$ is linear, multiplicative, preserves the involution, and is norm-continuous,
 			\item for every $A\in\A$, the map $t\mapsto \alpha_t(A)$ is continuous in the norm topology.
 		\end{enumerate}
 	\end{definition}
 	
 	By duality, the time evolution of states is given by $(\alpha_t^*\omega)(A)=\omega(\alpha_t(A))$ for all $A\in\A$. In addition, since $\alpha_t$ are $^*$-automorphisms, they preserve every selfadjoint sector $\A_{sa}$, ensuring that real physical observables evolve into real physical observables.
 	
 	\smallskip
 	\paragraph{Conclusion.} \emph{The physics of a classical system can be rewritten algebraically as the pair $(\A,\alpha_t)$ where $\A$ is a commutative unital $C^*$-algebra as in Definition \ref{def:unital-c-star-algebra} and $\alpha_t$ is a time evolution on it, as in Definition \ref{def:time-evolution}. To fix the initial conditions, it is given a state $\omega \in \mathcal{S}(\A)$, as in Definition \ref{def:state}.}
	 
\subsection*{Characters of a $C^*$-algebra and Gelfand-Naimark Theorem}
In the previous discussion, we have concluded that Hamiltonian mechanics can be rewritten in algebraic terms. In this section, we discuss in detail the methods required to derive a state space from a $C^*$-algebra. We demonstrate how the definition of a commutative $C^*$-algebra can be exploited to obtain a topological space through the Gelfand-Naimark theorem, which will act as the state space of the system. We achieve this using a class of multiplicative linear functionals called characters. These characters coincide with the pure states and their space will coincide with the actual state space.

\begin{definition}[Character of a commutative $C^*$-Algebra]\label{def:character}
	Let $\A$ be a commutative $C^*$-algebra, as in Definition \ref{def:c-star-algebra}. A \emph{character} $\omega$ of $\A$ is a nonzero linear functional $\omega: \A \to \C$ that preserves algebraic multiplication, meaning
	$$
	\omega(AB) = \omega(A)\omega(B),
	$$
	for all $A, B \in \A$. The \emph{Gelfand spectrum} $\sigma(\A)$ is the set of all characters on $\A$.
\end{definition}

Characters map observables to complex numbers and naturally belong to the dual space of the algebra. To analyse their geometric properties, we must equip this dual space with a metric structure.

\begin{definition}[Dual Operator Norm]\label{def:dual-norm}
	Let $\A$ be a $C^*$-algebra as in Definition \ref{def:c-star-algebra}. The dual space $\A^*$ is a Banach space as in Definition \ref{def:banach-space} equipped with the dual operator norm, defined for any functional $\alpha \in \A^*$ as
	$$
	\|\alpha\| = \sup_{A \in \A} \{\  |\alpha(A)| \ | \ \|A\| \le 1 \}.
	$$
\end{definition}

\noindent This definition is needed to understand the boundedness and continuity of characters that make the Gelfand's spectrum a unit ball.

\begin{lemma}[Boundedness of characters]\label{lemma:bound_char}
    Let $\A$ be a commutative $C^*$-algebra, as in Definition \ref{def:c-star-algebra}, and its Gelfand spectrum $\sigma(\A)$ per Definition \ref{def:character}. We have $\sigma(\A) \subseteq \A^*$ and for any $\omega \in \sigma(A)$, $\|\omega\| = 1$. Furthermore, if $\A$ is unital then $\omega(\Id_\A) = 1$.
\end{lemma}

For a complete proof of this lemma, we refer to \cite[Lemma C.9]{landsmanFoundationsQuantumTheory2017}. So far, characters have been defined only abstractly, but, in fact, they are exactly what we earlier called pure states. To establish this, we need the following result proved in \cite[Proposition 2.3.27]{bratteliOperatorAlgebrasQuantum}.

\begin{proposition}[Equivalence of Pure States and Characters]\label{prop:equivalence_states_characters}
	Let $\omega$ be a nonzero linear functional over the commutative $C^*$-algebra $\A$. Then $\omega$ is a pure state if and only if $\omega \in \sigma(\A)$.
\end{proposition}

Now, we need a topological structure for the spectrum to behave as a classical state space. The uniform topology induced by the dual norm is inadequate because in any infinite dimensional Banach space, the closed unit ball is never compact with respect to the uniform norm, see Riesz Lemma \cite[Theorem II.4]{reedMethodsModernMathematical2003}. We therefore introduce a coarser topology designed to recover this required compactness.

\begin{definition}[Weak$^*$ Topology]\label{def:weak_topologies}
	Let $\A^*$ be the topological dual space of a $C^*$-algebra $\A$, as in Definition \ref{def:c-star-algebra}. The \emph{weak$^*$ topology} on $\A^*$ is the coarsest topology such that for every fixed $A \in \A$ the evaluation map $\hat{A} : \A^* \to \C$ defined by $\hat{A}(\alpha) = \alpha(A)$ is continuous $\forall \alpha \in \A^*$.
\end{definition}

This topology controls pointwise convergence. A sequence of functionals converges if and only if the expectation value of every individual observable converges. By relaxing the strict requirement of uniform convergence across the entire algebra, this topology restores the geometric compactness essential for a well defined physical phase space.

\begin{proposition}[Weak$^*$ Compactness of the Spectrum]\label{prop:spectrum-compactness}
	Let $\A$ be a unital commutative $C^*$-algebra as in Definition \ref{def:c-star-algebra}. The Gelfand spectrum $\sigma(\A)$, defined in Definition \ref{def:character}, is a compact Hausdorff space with respect to the weak$^*$ topology.
\end{proposition}
\noindent The complete proof of this result can be found in \cite[Lemma C.10]{landsmanFoundationsQuantumTheory2017}.
	 
	 \noindent Now, we present the central result of this chapter. This theorem allows us to completely reconstruct the topological state space of a physical classical system directly from its abstract commutative algebra, establishing a correspondence between algebraic structures and geometry.
	 
	 \begin{theorem}[Gelfand-Naimark Representation Theorem]\label{thm:gelfand}
	 	Let $\A$ be a commutative $C^*$-algebra and $X$ be the Gelfand spectrum of $\A$, as in Definition \ref{def:character}, equipped with the weak$^*$ topology induced by $\A^*$ as in Definition \ref{def:weak_topologies}. Then $\A$ is isomorphic to the algebra $C_0(X)$ of continuous functions over $X$ which vanish at infinity. 
	 \end{theorem}
	 
	 \noindent A complete proof of this theorem is provided in \cite[Theorem 2.1.11]{bratteliOperatorAlgebrasQuantum}. As a special case we state this corollary that is also proved in \cite[Theorem 2.1.11A]{bratteliOperatorAlgebrasQuantum}.
	 \begin{corollary}\label{cor:gelfand}
	 	Given the hypothesis of Theorem \ref{thm:gelfand}, $X$ is a locally compact Hausdorff space which is compact if, and only if, $\A$ is unital, as in Definition \ref{def:unital-c-star-algebra}. Thus, we have $\A \cong C^0(X)$, \textit{i.e.}, the space of continuous functions over $X$.
	 \end{corollary}
	 
The Gelfand-Naimark Theorem \ref{thm:gelfand} establishes the duality necessary to reconstruct the state space from any commutative $C^*$-algebra. It yields a locally compact Hausdorff space $X$ that matches the Gelfand spectrum formed by the characters of the algebra. In addition, this spectrum is isomorphic to the classical continuous algebra, seen in Definition \ref{def:gel-alg}.

The topology of $X$ directly reflects the algebraic structure of $\A$.
As formalized in Proposition \ref{prop:spectrum-compactness}, the existence of
the identity element ensures that the Gelfand spectrum forms a weak$^*$ closed
subset of the dual unit ball, which guarantees its strict compactness via the
Banach-Alaoglu Theorem \cite[Theorem IV.21]{reedMethodsModernMathematical2003}. Consequently, the Gelfand isomorphism maps the abstract
algebra to the space $C_0(X)$ of continuous functions vanishing at infinity,
which identically reduces to the full space of continuous functions $C^0(X)$
in this compact setting, as stated in Corollary \ref{cor:gelfand}.


\subsection*{Reconstruction of the Symplectic Manifold}

The Gelfand isomorphism identifies the topological spectrum $X$ purely as a Hausdorff space, discarding all differentiable and symplectic information. To recover the phase space of classical mechanics, the $C^*$-algebraic framework must be augmented with additional algebraic structures. These structures are captured by the following two definitions. The Poisson algebra formalizes the Lie bracket of observables onto the $C^*$-algebra itself, while the smooth dense subalgebra supplies the required differentiability in a subset of $\A$. The compatibility between them, enforced by the Leibniz rule, ensures that the resulting geometry is coherent. However, the choice of such a subalgebra is not inherent to $\A$ itself and constitutes the central external input of this reconstruction procedure.

We begin by defining the main structure that allows us to recover the Lie algebra characteristic of classical mechanics.

\begin{definition}[Poisson Algebra]\label{def:poi-alg}
	A Poisson algebra is an associative commutative algebra \(\mathfrak P\), as in Definition \ref{def:algebra}, equipped with a Lie bracket \(\{\cdot,\cdot\}\), per Definition \ref{def:lie-algebra}, acting as a derivation with respect to the associative product. Namely it satisfies the Leibniz rule
	\[
	\{f, lh\} = \{f, l\}h + l\{f, h\} \qquad \forall f, h, l \in \mathfrak P.
	\]
\end{definition}

In addition, to recover the smooth structure of the phase space, we need an isomorphism from, at least, a subalgebra of $\A$ to a particular $C^\infty_0(M)$ where $M$ is a smooth manifold. These strong requirements are coded inside the following definition.

\begin{definition}[Smooth Dense Poisson Subalgebra]\label{def:smooth-subalgebra}
	Let \(\A\) be a commutative unital C*-algebra as in Definition \ref{def:unital-c-star-algebra} with Gelfand spectrum \(X\) as in Definition \ref{def:character}. A \emph{smooth dense Poisson subalgebra} is a Poisson subalgebra $(\A^\infty, \{\cdot,\cdot\})$, seen in Definition \ref{def:poi-alg}, with \(\A^\infty \subset \A\) satisfying
	\begin{enumerate}
		\item  \(\A^\infty\) is a selfadjoint subset as in Definition \ref{def:selfadj-subset},
		\item \(\A^\infty\) is dense in \(\A\) with respect to the uniform norm,
		\item there exists a smooth manifold \(M\) such that \(\A^\infty \cong C^\infty_0(M)\) as algebras.
	\end{enumerate}
\end{definition}

 We can now ask if the Poisson and the smooth structure defined on $\A^\infty$ can be transported and extended to the whole space of characters $\sigma(\A) \equiv X$. The first answer is negative due to the behaviour of the Poisson brackets, as explained in the next Remark \ref{rmk:obstruction-poi-strc}.

\begin{remark}[Topological Obstruction to Poisson Structure]\label{rmk:obstruction-poi-strc}
	The Poisson bracket on the Poisson subalgebra, as seen in Definition \ref{def:poi-alg}, acts as a first-order differential operator. This makes it generally unbounded and discontinuous with respect to the uniform supremum norm of the C*-algebra $\A$. This obstruction restricts the Poisson algebra strictly to the incomplete dense subalgebra $\A^\infty$, sacrificing Banach completeness to preserve the Leibniz rule.
\end{remark}

On the other hand, when we examine the smooth structure instead, we find that there is a technical way to extend it. In fact, if there exists a smooth structure $M$, as in Definition \ref{def:smooth-subalgebra}, and it is compact, we can prove the following result.

\begin{proposition}
	Let \(\A\) be a commutative unital C*-algebra as in Definition \ref{def:unital-c-star-algebra} with Gelfand spectrum \(X\) as in Definition \ref{def:character}. Let \((\A^\infty, \{\cdot,\cdot\})\) be the smooth dense Poisson subalgebra of $\A$, as in Definition \ref{def:smooth-subalgebra}, with \(\A^\infty \cong C_0^\infty(M)\). If \(M\) is compact, then \(X\) is a smooth manifold.
\end{proposition}

\begin{proof}
Since $M$ is compact, $C_0^\infty(M) = C^\infty(M)$ contains the constant functions. By the isomorphism of Definition \ref{def:smooth-subalgebra}, it descends that $\Id_{\A} \in \A^\infty$. In particular $\A^\infty$ is a unital subalgebra of $\A$. Let $\sigma(\A^\infty)$ denote the set of linear multiplicative non-trivial functionals on $\A^\infty$, and define an evaluation map as
\[
\mathrm{ev} : X \to \sigma(\A^\infty), \qquad
\mathrm{ev}(x) = \mathrm{ev}_x := x|_{\A^\infty}, \qquad
\mathrm{ev}_x(f) = f(x) \quad \forall f \in \A^\infty.
\]
That is well defined and continuous by \cite[Construction p.~92]{landsmanFoundationsQuantumTheory2017} being dual to the inclusion $\mathfrak{A}^\infty \hookrightarrow \mathfrak{A}$.

\noindent Since $\A^\infty$is dense in $\A$ and multiplication is jointly continuous, two continuous maps out of $\A$ that agree on $\A^\infty$ must coincide, and likewise two continuous maps out of $\A \times \A$ that agree on $\A^\infty \times \A^\infty$ must coincide, via \cite[Exercise 13, \S 18]{munkresTopology2018}.

\begin{description}[wide=0pt, leftmargin=*, font=\normalfont\itshape]
	\item[Injectivity.] If $x_1|_{\A^\infty} = x_2|_{\A^\infty}$, the above result yields $x_1 = x_2$ on $\A$.
	\item[Surjectivity.] Under $\A^\infty \cong C_0^\infty(M)$, every $\omega \in \sigma(\A^\infty)$ induces an evaluation by duality $\omega(f) = f(p)$ $\forall f \in C_0^\infty(M)$ for a particular $p \in M$. Since all $f \in \A^\infty$ are bounded by the supremum norm inherited from $\A$, $\omega$ is bounded. Thus the B.L.T.\ Theorem \cite[Theorem I.7]{reedMethodsModernMathematical2003} extends the functionals uniquely to a bounded linear $\tilde\omega \in \A^*$. Since the maps $(A,B) \mapsto \tilde\omega(AB)$ and $(A,B) \mapsto \tilde\omega(A)\tilde\omega(B)$, $\forall A,B \in \A$, agree on the dense subset $\A^\infty \times \A^\infty$, the same above fact yields $\tilde\omega(AB) = \tilde\omega(A)\tilde\omega(B)$. It descends that $\tilde\omega \in \sigma(\A)$ and $\mathrm{ev}$ is surjective.
\end{description}

\noindent It descends that $\mathrm{ev}$ is a continuous bijection. It remains to identify $M$ with $X$. Under $\A^\infty \cong C^\infty(M)$, the evaluation map
\[
\mathrm{ev}^M : M \to \sigma(\A^\infty), \qquad
p \mapsto \mathrm{ev}^M_p, \qquad
\mathrm{ev}^M_p(f) = f(p) \quad \forall f \in \A^\infty,
\]
is a bijection \cite[Proposition~C.21]{landsmanFoundationsQuantumTheory2017}. Continuity follows since the evaluation $\mathrm{ev}^M_p(f) = f(p)$ acts on smooth, hence continuous, functions $f \in \A^\infty$. Since $M$ is compact and $\sigma(\A^\infty)$ is Hausdorff, $\mathrm{ev}^M$ is a homeomorphism by \cite[Theorem~26.6]{munkresTopology2018}:
$$M \cong \sigma(\A^\infty).$$
\noindent Since $\A$ is unital, $X$ is compact by Gelfand Theorem Corollary \ref{cor:gelfand}. Also, since $\sigma(\A^\infty)$ is Hausdorff, the continuous bijection $\mathrm{ev}$ is likewise a homeomorphism via \cite[Theorem 26.6]{munkresTopology2018}:
$$X \cong \sigma(\A^\infty).$$
Composing the two homeomorphisms yields $X \cong M$. Now, transporting the smooth atlas of $M$ along this homeomorphism equips $X$ with a smooth structure, making it a smooth manifold diffeomorphic to $M$.
\end{proof}

\noindent Having discussed the limitations and possible extension of the smooth subalgebra $\A^\infty$, we can collect the components identified so far into a single, self-contained structure.

\begin{definition}[Classical Algebraic System]\label{def:cas}
	A classical mechanical system is completely characterized by the triplet $(\A, \A^\infty, \{\cdot,\cdot\})$, where
	\begin{enumerate}
		\item $\A \cong C^0(X)$ is a commutative unital C*-algebra, as in Definition \ref{def:unital-c-star-algebra}, determining the topological state space $X$, its Gelfand spectrum.
		\item $(\A^\infty \subset \A, \{\cdot, \cdot\})$ is a smooth dense Poisson subalgebra, as in Definition \ref{def:smooth-subalgebra}, inducing the differential and dynamical structure.
	\end{enumerate}
\end{definition}

The last step is to show that this is sufficient to recover the complete structure of a  symplectic manifold. This can be done invoking the representation theory of Poisson algebras. The complete explanation of this theory is outside the scope of this work, however we can refer the reader to the representation theorem presented in \cite[Theorem 2.6.7]{landsmanMathematicalTopicsClassical1998}. The key idea is to use the Poisson bracket defined on the smooth dense subalgebra to endow the smooth manifold $M$ with a dynamical structure. The theorem establishes that irreducible representations of the Poisson subalgebra, presented in Definition \ref{def:smooth-subalgebra}, correspond to realizations of the algebra on spaces where the Poisson brackets become strictly non-degenerate. Geometrically, this implies evaluating the global Poisson structure of $M$ on invariant submanifolds that inherit a well-defined symplectic form. By restricting the dynamics to these irreducible geometric domains, one exactly recovers the full symplectic phase space formulation of Hamiltonian mechanics without additional postulates.

\subsection*{Physical Observables and Dynamics}\label{sec:interpret-phys}
Before ending this chapter, we clarify the physical interpretation of the algebras introduced so far. We start with a unital commutative $C^*$-algebra $\A$, per Definition \ref{def:c-star-algebra}. Using Gelfand Theorem \ref{thm:gelfand}, one can reconstruct the topological space $X$. The selfadjoint sector $\A_{sa}$, already seen in Definition \ref{def:obs-alg}, isolates the measurable physical quantities. The smooth dense Poisson subalgebra $\A^\infty$ of Definition \ref{def:smooth-subalgebra} is required to formulate Hamiltonian mechanics. These topological, algebraic, and differential requirements combine as follows.
\begin{itemize}
	\item The continuous measurable observables form the selfadjoint sector $\A_{sa} = \{f \in \A \mid f^* = f\}$. This space suffices to define states, probability measures, and continuous symmetries, but lacks a Poisson bracket and so cannot generate differential dynamics.
	\item Hamiltonian mechanics further requires smoothness to formulate the equations of motion, so the Hamiltonian $H$ must belong to the smooth selfadjoint sector $\A^\infty_{sa} := \A^\infty \cap \A_{sa}$. This guarantees that the Hamiltonian vector field $X_H = -\{H, \cdot\}$ generates a well-defined, strictly real one-parameter group of $^*$-automorphisms $\alpha_t : \A^\infty \to \A^\infty$.
\end{itemize}
A non-selfadjoint or non-smooth Hamiltonian would produce complex flows or exit the domain of the Poisson bracket, yielding unphysical dynamics. Restricting dynamical generators to the selfadjoint part of $\A^\infty$, preserves both the real structure of physical observables and the symplectic geometry of the underlying phase space.

\newpage
\section{Example I: Classical Spin System}\label{sec:1-example}
In this section, we apply the algebraic tools developed in the first chapter to the classical spin system. This example provides an interesting test case for our framework, since it is a non-trivial Poisson algebra whose reduction yields a compact curved manifold, the sphere. This allows every abstract construction introduced so far, from the Gelfand spectrum to the Poisson bracket, to be verified explicitly on a concrete and geometrically intuitive example. Working through this case in detail also illustrates how symmetries of a physical system translate, via Noether's Theorem, into conserved quantities within the same algebraic language.

\subsection{The Poisson Algebra and Gelfand's Theorem}

Consider the polynomial algebra \(\A \), as in Definition \ref{def:algebra}, generated by three elements \(x_1, x_2, x_3\), equipped with the \(\mathfrak{so}(3)\)-Poisson bracket
\[
\{x_i, x_j\} = \epsilon_{ijk} x_k, \qquad i,j,k \in \{1,2,3\},
\]
where \(\epsilon_{ijk}\) is the Levi-Civita symbol. This bracket is skew-symmetric and satisfies the Jacobi identity by virtue of the Lie algebra structure. A direct computation reveals the existence of a Casimir function
\[
C = x_1^2 + x_2^2 + x_3^2, \qquad \{C, A\} = 0 \quad \forall A \in \A.
\]

Since \(\A\) is abstract, we must first extract a space from it. In order to fix \(C=1\), we identify \(A \sim A'\) whenever \(A-A' = Q(C-1)\) for some \(A, A', Q \in \A\). Call the resulting quotient \(\A_1\). This identification respects the bracket because \(C\) is a Casimir operator that acts as a scalar through the Poisson brackets. It descends that \(\{A,B\}-\{A',B\} = (C-1)\{Q,B\}\) for $B \in \A$, which again vanishes in the quotient. It descends that \(\A_1\) inherits a well-defined Poisson brackets.

The points of \(\A_1\) are the multiplicative linear functionals \(\omega:\A_1\to\R\) in the dual space \(\A_1^*\). We call their collection \(X\). Since \(x_1,x_2,x_3\) generates \(\A_1\), each \(\omega\in X\) is fixed by the triple \((\omega(x_1),\omega(x_2),\omega(x_3))\in\R^3\). Multiplicativity and \(C=1\) in \(\A_1\) force this triple onto \(\mathbb S^2\). In contrast, every point of \(\mathbb S^2\) defines such a functional by evaluation. Hence,
\[
X \cong \mathbb{S}^2.
\]
Completing \(\A_1\) in the uniform norm on \(X\) yields a commutative $C^*$-algebra isomorphic to \(C^0(\mathbb S^2)\), consistent with Gelfand's Theorem \ref{thm:gelfand}. Moreover, the brackets descend with it, making \(\mathbb S^2\) a natural candidate for a symplectic manifold, seen in Definition \ref{def:sympl-man}.

\subsection{Construction of the Symplectic Manifold}
We now equip $\mathbb{S}^2$ with coordinates and explicitly construct the symplectic form. In cartesian coordinates, the symplectic form in $\R^3$ is given by 

$$\omega =  x_1 dx_2 \wedge dx_3 + x_2 dx_3 \wedge dx_1 + x_3 dx_1 \wedge dx_2. $$

\noindent Now we restrict the form to $\mathbb{S}^2$ using the standard spherical coordinates
\[
	x_1 = \sin\theta \cos\phi, \qquad x_2 = \sin\theta \sin\phi, \qquad x_3 = \cos\theta,
\]
where \(\theta \in [0,\pi] \) and \(\phi \in [0, 2\pi)\). Now we must compute the bracket \(\{\theta, \phi\}\). Using the chain rule and computing the partial derivatives we get%From \(\{x_1, x_2\} = x_3 = \cos\theta\), we have
%\[
%	\{x_1, x_2\} = \frac{\partial x_1}{\partial \theta}\frac{\partial x_2}{\partial \phi}\%{\theta,\phi\} + \frac{\partial x_1}{\partial \phi}\frac{\partial x_2}{\partial \theta}\%{\phi,\theta\}.
%\]
%Substituting the partial derivatives yields
\[
	\{x_1, x_2\} = \sin\theta \cos\theta \{\theta,\phi\}.
\]
Equating this to \(x_3 = \cos\theta\) yields
\begin{equation}
	\boxed{\{\theta, \phi\} = \frac{1}{\sin\theta}.}\label{eq:poi-sphere}
\end{equation}
\noindent This formula is valid away from the poles. Its extension to \(\sin\theta=0\) is given by the Cartesian expression, ensuring global non-degeneracy.
Given that the manifold is two-dimensional, to compute the symplectic form in spherical coordinates, we invert the commutator found in Equation \eqref{eq:poi-sphere} to obtain
\begin{equation}
	\boxed{\omega = \sin\theta \, d\theta \wedge d\phi.} \label{eq:omega-sphere}
\end{equation}
\noindent This form is closed \(d\omega = 0\) and non-degenerate, thus \((\mathbb{S}^2, \omega)\) is a symplectic manifold as in Definition \ref{def:sympl-man}.

\subsection{Hamiltonian Action and the Moment Map}

The group \(G = SO(3)\) acts on \(\mathbb{S}^2\) by rotations. For each \(\xi \in \mathfrak{so}(3) \cong \mathbb{R}^3\), the infinitesimal generator, seen in Definition \ref{def:inf-gen}, is given by
\[
	X_\xi(x) = x \times \xi,
\]
where \(\times\) denotes the standard cross product. We seek a moment map \(\mu: \mathbb{S}^2 \to \mathfrak{so}(3)^*\), as in Definition \ref{def:mom-map}, satisfying the relation
\[
	d\langle \mu(x), \xi \rangle = \iota_{X_\xi} \omega, \qquad \forall \xi \in \mathfrak{so}(3).
\]
Identifying \(\mathfrak{so}(3)^* \cong \mathbb{R}^3\) via the dot product, we test the candidate \(\mu(x) = -x\). Given \(\xi = e_1 = (1,0,0)\), then \(X_{e_1} = (0, x_3, -x_2)\) and contracting \(\omega\) with \(X_{e_1}\) yields
\begin{align*}
	\iota_{X_{e_1}} \omega 
	&= \iota_{(0,x_3,-x_2)} \left( x_1 dx_2 \wedge dx_3 + x_2 dx_3 \wedge dx_1 + x_3 dx_1 \wedge dx_2 \right) \nonumber \\
	&= x_1(x_3 dx_3 + x_2 dx_2) - (x_2^2+x_3^2) dx_1.
\end{align*}

\noindent Using the constraint \(x_2 dx_2 + x_3 dx_3 = -x_1 dx_1\), given by the restriction on the sphere, we get
\[
	\iota_{X_{e_1}} \omega = -dx_1.
\]
\noindent Thus, \(d\langle \mu, e_1 \rangle = d(-x_1)\), which implies
\begin{equation}
	\boxed{\mu(x) = -x.}\label{eq:mom-map}
\end{equation}
\noindent Physically, this is the negative of the position vector, which corresponds to the angular momentum vector, up to a sign convention.

\subsection{Noether's Theorem in Practice}
The moment map constructed above is the algebraic counterpart of a conserved quantity associated with a continuous symmetry, in the sense of Noether's Theorem \ref{thm:noether}. We now illustrate this correspondence through the example of the precession of a classical system that is governed by the same underlying symplectic structure on the sphere.
\medskip
Consider the Hamiltonian
\[
	H = x_3,
\]
which is invariant under rotations about the \(z\)-axis. The symmetry generator is \(\xi = e_3 = (0,0,1)\). The associated Noether invariant is
\[
	\langle \mu, \xi \rangle = \langle -x, e_3 \rangle = -x_3.
\]
As calculated in Equation \eqref{eq:mom-map}, Noether's theorem guarantees \(x_3\) is conserved. In fact, the equations of motion are
\begin{align*}
	\dot{x}_1 &= \{x_1, H\} = -x_2, \\
	\dot{x}_2 &= \{x_2, H\} = x_1, \\
	\dot{x}_3 &= \{x_3, H\} = 0.
\end{align*}
Solving yields a uniform precession around the \(z\)-axis:
\[
	x_1(t) + i x_2(t) = e^{it} \big(x_1(0) + i x_2(0)\big), \qquad x_3(t) = x_3(0).
\]

\printbibliography
\end{document}
